% ============================================================
% Section 1 — Introduction (v4, 2026-09-20: 压缩 + 第一页右栏机制图占位 + 去掉安置贡献)
% v3 (2026-09-16/18) 被压缩/删除的部分以注释保留在文末; v2 原句见 git 历史。
% 数据: HANDOFF §22 (K=8 修正), §29/§30 (查询级信用 seed 1), §31-§37 (开放续训), §33 (训练后 K=8)
% ============================================================

\section{Introduction}
\label{sec:intro}

  Multi-hop question answering requires an agent to acquire evidence
  sequentially: each retrieval exposes an intermediate entity that the next
  query depends on. Reinforcement learning has become the standard way to
  train such search agents---the policy interleaves reasoning, search, and
  answering over several turns, optimized by a critic-free policy
  gradient whose baseline is the mean reward of the $K$ rollouts drawn
  for the same question---group-relative advantage estimation, as in
  GRPO and leave-one-out
  REINFORCE~\citep{shao2024deepseekmath,ahmadian2024rloo,kool2019buy}---
  against a terminal reward for answer
  correctness~\citep{jin2025searchr1,chen2025research,song2025r1searcher,
  sun2026zerosearch}. This recipe improves aggregate accuracy, but its gains
  concentrate on shallow questions; on the deepest chains, where sequential
  evidence acquisition matters most, improvements remain modest.

  The standard diagnosis is temporal~\citep{zheng-etal-2025-stepsearch}:
  rewards are sparse along the trajectory, so intermediate progress goes
  uncredited. We argue for a
  different, \emph{group-level} diagnosis. Under group-relative advantage
  estimation, the policy gradient exists only where the $K$ rollouts of
  the same question differ in return; identical returns cancel exactly.
  On hard questions the terminal reward routinely fails to differ: every
  rollout is wrong, or every rollout is right. We call such a group
  \emph{answer-tied}. In deep search this is the dominant case: at the
  SFT initialization, 56--71\% of rollout groups are answer-tied at every
  depth. Depth changes which kind of tie dominates. At two hops, 28\% of
  groups are tied because every rollout is right; at four hops that share
  falls to 4\%, and 57\% of all groups are tied because every rollout is
  wrong (Fig.~\ref{fig:teaser}). In an
  answer-tied group the terminal reward is the same for all $K$ rollouts
  and contributes nothing to any advantage: a rollout that retrieved half
  of the gold evidence chain and one that retrieved none receive identical
  reward. The ordering that would separate them---how much of the
  evidence each rollout reached---is the signal this paper adds. What
  remains in an all-wrong group is the cost terms---search cost, token
  cost, format penalties---so the rollout that searched least ranks
  highest. On the questions where the agent most needs to learn to keep
  searching, the training signal rewards stopping early.

% ---- Fig 1 占位(单栏, 第一页右栏顶; 用户 2026-09-20: 先写好位置) ----
% 内容建议: 一个 4-hop 全错组的 K 条 rollout, 左列 outcome-only 优势(按搜索次数排),
% 右列查询级信用(按证据到达量排); 或 tie 账本三柱。成图后换 \includegraphics。
\begin{figure}[t]
\centering
\fbox{\parbox{0.95\columnwidth}{\centering\vspace{7em}
[Figure~1 placeholder, single column: an answer-tied group under the
answer reward (identical reward; the cost terms rank the rollout that
searched least first) and under query-level credit (the query that
reached new evidence is credited).]\vspace{7em}}}
\caption{Answer ties and their repair. [Placeholder.]}
\label{fig:teaser}
\end{figure}

  This failure mode has been named in mathematical RLVR as
  \emph{advantage collapse}~\citep{he2026collapse}, and existing
  mitigations discard uniform-outcome
  groups~\citep{yu2025dapo,greso2025,xu2026downsampling}, inject virtual
  reward samples~\citep{he2026collapse}, or re-baseline the comparison
  by trajectory structure~\citep{zhu2025grpo,Ji2026tree}. None gives a
  tied group an ordering with content; and left to chance, whatever
  residual signal occupies an all-failure group can come to govern the
  learned policy~\citep{wang2026darkroom}. Standard multi-hop training
  sets already contain the information that breaks these ties: gold
  supporting documents~\citep{trived2022musique,yang2018hotpotqa}.

  Where this information enters the estimator decides what the policy
  learns from it. Added to the trajectory return as an evidence-coverage
  term, it reorders whole trajectories: one advantage per rollout,
  broadcast to every turn, so the search step that reached new evidence
  and the answer step that followed receive the same sign; the policy
  learns \emph{when to answer}, a habit that pays inside the training
  environment and costs outside it (\S\ref{sec:mechanism}). We therefore
  assign the credit at the query step. At each search turn we sample $M$
  additional queries from the current policy at the same state, retrieve
  for each, and credit the executed query by its gold-evidence gain
  relative to the mean gain of the samples: a state-conditional baseline
  estimated from the policy's own action distribution, as in the vine
  estimator of \citet{schulman2015trpo}. The credit needs the dataset's
  gold annotations and the retriever, nothing else, and is used only
  during training. In an answer-tied group it is non-zero wherever the
  executed query and its alternatives reach different evidence, which is
  where the tie is broken.

  Our experiments on MuSiQue with a Qwen3.5-9B policy yield three
  findings. \textbf{First}, query-level credit improves deep search in
  the training environment: 51.4 exact match on the closed-pool
  development set against 47.6 for outcome-only training ($z=4.7$), with
  4-hop EM 40.0 against 31.1 ($z=4.1$), higher final evidence coverage,
  and 0.6 fewer search calls per question. \textbf{Second}, a
  group-level audit locates the gain: adding coverage expands the share
  of 4-hop groups with an informative ordering from 38\% to 89\%, and on
  the trained policy the share of 4-hop groups in which every rollout is
  wrong falls from 42.9\% to 37.3\% ($z=2.7$). \textbf{Third}, the gains
  carry to an open 21M-passage corpus: after 50 episodes against the
  open retriever the policy reaches 28.6 EM on MuSiQue and 48.6 on
  HotpotQA, above every search-RL system reproduced under the same
  retriever (ReSearch 22.1, Search-R1 45.0), with the credit adding
  $+2.7$ on 2WikiMultihopQA over the same continued training without it
  ($z=3.7$).

  Our contributions:
  \begin{itemize}\setlength\itemsep{2pt}
  \item \textbf{A group-level diagnosis.} Answer-tied groups dominate
    deep-search RL (56--71\% of groups), shift with depth from all-right
    ties to all-wrong ties (57\% of 4-hop groups), and carry a
    cost-driven residual ordering that favors the rollout that searched
    least (\S\ref{sec:method}).
  \item \textbf{Query-level evidence credit.} A turn-level credit for
    search queries, computed against a state-conditional baseline from
    the policy's own samples and the dataset's gold annotations
    (\S\ref{sec:method}, \S\ref{sec:mechanism}).
  \item \textbf{Deep-hop gains and the strongest open-corpus results.}
    $+3.8$ EM and $+8.9$ 4-hop EM over outcome-only training in the
    closed pool; on the open corpus, 28.6 EM on MuSiQue and 48.6 on
    HotpotQA after continued training against the retriever, above all
    six reproduced search-RL systems (\S\ref{sec:results}).
  \end{itemize}

% ---------------------------------------------------------------
% v3 原文(2026-09-16 至 09-18)中被压缩/删除的部分, 供回溯:
% 第 3 段(三派): "This failure mode has recently been named in mathematical
% RLVR---\emph{advantage collapse}---and existing mitigations fall into three
% families. Discarding methods drop or skip uniform-outcome groups by dynamic
% sampling or filtering, recovering compute but abandoning exactly the hardest
% questions. Synthetic-advantage methods inject virtual reward samples so that
% collapsed groups still produce a learning signal, but the signal is
% content-free: the ordering inside a tied group remains arbitrary. Structural
% methods re-baseline the comparison itself, stratifying advantages by
% trajectory structure or comparing along shared tree prefixes; this repairs
% cross-question fairness but leaves a fully tied group as uninformative as
% before. Leaving such groups to chance is not neutral: in GRPO training,
% whatever residual signal occupies an all-failure group can come to govern
% the learned policy outright. Tied groups should be neither thrown away nor
% filled with noise; what they need is an ordering that \emph{means}
% something."
% 第 4 段(方法)原有: "This placement teaches the policy \emph{when to
% answer}---withhold until the evidence set is complete---a habit that pays
% inside the training environment and costs outside it (\S mechanism,
% \S results-open)." / "Answer turns receive the outcome reward alone."
% 发现一原有: "against 47.6 for outcome-only training and 47.5 for the
% trajectory-level coverage term from the same initialization and seed
% ($z=4.7$ and $4.9$)"; 发现二原有: "at the training start, 84\% of 4-hop
% answer-tied groups carry coverage variation the answer reward cannot see";
% 发现三原有: "Without any training against it, the query-level arm scores
% $+3.4$ EM over the trajectory-level term on 2WikiMultihopQA ($z=4.5$) and
% $+4.0$ at four hops on MuSiQue ($z=2.4$);"
% 贡献一原有尾句: "---an account of deep-hop stagnation that temporal sparsity
% does not capture, complementing advantage-collapse analyses in mathematical
% RLVR~\citep{yu2025dapo,he2026collapse}"; 贡献二原有: "with the
% trajectory-level coverage term as the controlled contrast that isolates
% where the signal must enter".
% 贡献四(已撤, 安置研究进附录): "\textbf{A signal-placement study.} The same
% evidence signal helps in training and offline audit but not as runtime
% control or injected observation (\S placement)."
% 发现四(注释版, 已撤): "\textbf{Fourth}, the same signal that helps as a
% training reward and as an offline audit yields no reliable gain as a runtime
% controller or as an injected observation: evidence supervision is
% placement-sensitive."
% ---------------------------------------------------------------
